\chapter{Related Work}
\label{chap:related}

This chapter surveys tools and research that overlap with one or more
aspects of the Pseudo++ toolchain: general-purpose online execution
platforms, pseudocode-specific interpreters, visual programming
environments, and the use of transpilation in educational settings.
Each category is assessed along the dimensions relevant to the Romanian
baccalaureate use case, and Section~\ref{sec:comparison} closes with a
direct feature comparison.

% ─────────────────────────────────────────────────────────────────────────────
\section{Online Code Execution Platforms}
\label{sec:online-editors}

The most commonly used tools in Romanian high-school computer science
classes are browser-based programming environments for the two target exam languages,
C++ and Pascal.  Platforms such as Replit \cite{replit} and OnlineGDB
\cite{onlinegdb} provide a full edit-compile-run loop in the browser
without any local installation, which is practical in lab settings where
students lack administrator rights.  Both support C++ and Pascal, and
OnlineGDB offers a graphical step-through debugger for C++.

The fundamental limitation shared by all such platforms is that the
input must be a real programming language.  A student who formulates a
solution in the baccalaureate pseudocode dialect must first translate it
manually into C++ or Pascal before it can be compiled.  This translation
step is an additional cognitive burden, and the compiled program does not
directly confirm the correctness of the pseudocode --- only of its
translation.  Neither platform provides any assistance with the
translation itself.

Python Tutor \cite{pythontutor}, developed by Guo at MIT, occupies a
complementary position.  It is a step-through visualiser for Python,
Java, C, and C++ that renders the heap, stack frames, and variable
bindings at each execution step.  Its variable-inspection panel --- a
live display of name, value, and type updated after every statement ---
is conceptually close to the variable inspector in Pseudo++'s web
editor.  Python Tutor's contribution is the visualisation model rather
than language coverage; like Replit it requires the student to work in a
general-purpose language, not in a pseudocode dialect.

% ─────────────────────────────────────────────────────────────────────────────
\section{Pseudocode-Specific Tools}
\label{sec:pseudocode-tools}

\textbf{PSeInt} \cite{psint} is the closest existing tool to Pseudo++.
Developed by Pablo Novara for the Spanish-speaking educational community,
it is a desktop application that interprets a pseudocode dialect with
Spanish keywords (\textit{Escribir}, \textit{Leer}, \textit{Si},
\textit{Mientras}) and provides a step-through debugger and a flowchart
view.  PSeInt has accumulated a large user base across Latin America and
has been actively developed for over two decades.

The structural differences from Pseudo++ are significant.  Being a
desktop application, PSeInt requires installation and is not accessible
from a browser.  It supports only its own Spanish-keyword dialect: it
does not recognise Unicode mathematical symbols, box-drawing characters,
or any of the Romanian baccalaureate keywords.  Critically, PSeInt has
no transpiler: there is no mechanism to convert a program to C++, Pascal,
or any other compilable language.

\textbf{RAPTOR} \cite{raptor} takes a fundamentally different approach to
pseudocode: its primary representation is a \emph{flowchart}, not text.
Students assemble programs from graphical boxes (assignment, decision,
loop, procedure call) connected by arrows.  The tool can display a
textual pseudocode rendering of the current flowchart, but this text is
a read-only projection of the graphical model, not an independently
editable source.  RAPTOR includes a visual debugger that highlights the
executing node in the flowchart.  It is widely used in introductory
programming courses at American universities \cite{raptor}, but its
visual-first design diverges substantially from the text-based pseudocode
taught in Romanian high schools, which students must also write and reason
about on paper during examinations.

\textbf{Flowgorithm} \cite{flowgorithm}, developed by Carlisle (also the
lead author of the RAPTOR paper), is a more recent flowchart tool with
a broader output capability.  Unlike RAPTOR it can export the constructed
flowchart to over a dozen programming languages, including C, C++, Java,
and Python.  This export feature is the closest existing analogue to
Pseudo++'s transpiler; the semantics-preserving translation from a
structured notation to compilable code serves the same pedagogical
purpose.  However, Flowgorithm's input is a graphical flowchart, not a
text-based pseudocode, so it cannot process the notation that appears on
Romanian exam papers.

\textbf{Pseudocod.ro} \cite{pseudocodro} is the only existing browser-based
tool that targets the Romanian pseudocode dialect directly.  It accepts
text input and executes it without requiring a local installation, which
places it closest to Pseudo++ among all surveyed tools.  However, its
accepted syntax diverges from the baccalaureate standard in several
respects, meaning that programs written in the notation used on exam
papers typically require manual adaptation before they run correctly.
The tool provides basic execution only: it has no step-through debugger,
no transpiler, no static analysis, and no loop-equivalence feature.

% ─────────────────────────────────────────────────────────────────────────────
\section{Visual and Block-Based Programming Environments}
\label{sec:educational-interpreters}

Block-based environments represent a distinct strand of educational
tooling, aimed primarily at learners who benefit from the scaffolding
that syntactically valid-by-construction programs provide.

\textbf{Scratch} \cite{scratch}, developed at the MIT Media Lab, is the
most widely deployed representative of this category.  Programs are
constructed by snapping coloured blocks together; there is no text to
type or syntax error to make.  Scratch targets an 8--16 age range and
its abstraction level --- event-driven sprites reacting to messages ---
is misaligned with the baccalaureate curriculum, which requires reasoning
about variables, arithmetic expressions, nested loops, and algorithms in
an imperative style.

\textbf{Blockly} \cite{blockly}, Google's open-source block editor
library, is lower-level than Scratch and can generate Python or
JavaScript source from its block graphs.  The block-to-code export path
is structurally analogous to transpilation: a high-level representation
is compiled to a lower-level, executable form.  Some educational
platforms use Blockly as a bridge to a real language, letting students
first construct a program visually and then inspect the generated
Python.  The key difference relative to Pseudo++ is the input
representation: Blockly's source is a visual block graph, not the
linear, text-based pseudocode notation that Romanian students write on
paper.

% ─────────────────────────────────────────────────────────────────────────────
\section{Comparison with the Present Work}
\label{sec:comparison}

Table~\ref{tab:comparison} summarises the capabilities of the surveyed
tools across seven dimensions.

\begin{table}[h]
\centering
\small
\setlength{\tabcolsep}{5pt}
\begin{tabular}{l c c c c c c c}
\hline
\textbf{Tool} &
\textbf{Executes} &
\textbf{Debugger} &
\textbf{Transpiles} &
\textbf{Linter} &
\begin{tabular}{@{}c@{}}\textbf{Loop-}\\\textbf{equiv.}\end{tabular} &
\begin{tabular}{@{}c@{}}\textbf{Browser-}\\\textbf{based}\end{tabular} &
\begin{tabular}{@{}c@{}}\textbf{Open}\\\textbf{source}\end{tabular} \\
\hline
Pseudocod.ro      & \checkmark & $\times$   & $\times$   & $\times$   & $\times$   & \checkmark & $\times$   \\
PSeInt            & \checkmark & \checkmark & $\times$   & $\times$   & $\times$   & $\times$   & \checkmark \\
RAPTOR            & \checkmark & \checkmark & $\times$   & $\times$   & $\times$   & $\times$   & $\times$   \\
Flowgorithm       & \checkmark & \checkmark & \checkmark & $\times$   & $\times$   & $\times$   & $\times$   \\
Replit            & \checkmark & partial    & $\times$   & $\times$   & $\times$   & \checkmark & $\times$   \\
Python Tutor      & \checkmark & \checkmark & $\times$   & $\times$   & $\times$   & \checkmark & \checkmark \\
Scratch           & \checkmark & $\times$   & $\times$   & $\times$   & $\times$   & \checkmark & \checkmark \\
Blockly           & $\times$   & $\times$   & \checkmark & $\times$   & $\times$   & \checkmark & \checkmark \\
\textbf{Pseudo++} & \checkmark & \checkmark & \checkmark & \checkmark & \checkmark & \checkmark & \checkmark \\
\hline
\end{tabular}
\caption{Feature comparison of educational programming tools.
  \textit{Linter}: static analysis pass that reports semantic warnings
  and errors independently of execution.
  \textit{Transpiles}: ability to produce compilable source in a
  standard language from the tool's primary input representation.
  \textit{Loop-equiv.}: interactive loop-equivalence transformation.
  \textit{Partial} for Replit's debugger reflects that step-through
  debugging is available for C++ but not for Pascal.}
\label{tab:comparison}
\end{table}

The table highlights a gap that none of the surveyed tools fills: the
combination of execution, step-through debugging, transpilation, linting,
and loop-equivalence transformations in a single browser-based,
open-source tool.  Flowgorithm comes closest --- execution, debugger, and
transpiler --- but it operates on graphical flowcharts rather than text
pseudocode and ships only as a desktop application.

Several qualitative differences complement the binary comparison.
PSeInt and RAPTOR require local installation, a practical barrier in
school lab environments.  Python Tutor's visualisation model is the
closest precedent for Pseudo++'s variable-inspection panel, but it
cannot process any form of pseudocode.  Replit and OnlineGDB expose
the student to the full syntactic constraints of C++ or Pascal before
any algorithm can be tested.
No existing tool natively executes the Romanian baccalaureate
pseudocode dialect: this remains the primary motivation for Pseudo++.
